\chapter{LIS - Longest Increasing Subsequence}

\section{Definizione del problema}

\subsection{Istanza}
Abbiamo una sequenze $\mathbf{X = \langle x_1, \dots, x_m \rangle (= X_m),\ \forall x_i \in X,\ x_i \in \Sigma}$

\subsection{Soluzione}
Una più lunga sottosequenza $Z$ di $X$ nella quale \uline{gli elementi appaiono in ordine crescente}.\\
($\Sigma$ munito di una relazione d'ordine crescente)

\subsection{Esempio}
{
    \large
    \begin{gather*}
        X = X_8 = \langle 6,4,7,3,5,9,11,8 \rangle \quad m=8 \\ \\ 
        \left.
        \begin{aligned}
        \langle 4, 7, 9, 11 \rangle \\
        \langle 6, 7, 9, 11 \rangle \\
        \langle 4, 5, 9, 11 \rangle \\
        \langle 3, 5, 9, 11 \rangle
        \end{aligned}
        \right\}
        \quad \text{Sono tutte LIS} \\ \\
        \langle 3,5,8\rangle  \text{ termina con $x_8 = 8$ ma non è una LIS perchè non ha lunghezza massima}
    \end{gather*}
}

\newpage
\section{Formulazione ricorsiva di una LIS}
\subsection{Istanza}
{
    \large
    $$
        X_i = \langle x_1, \dots, x_i \rangle
    $$
}

\subsection{Soluzione}
Una soluzione $LIS\ Z^{(i)}$ di $X_i$

\subsection{Caso base}
{
    \large

    \begin{itemize}
        \item $i=0\ Z^{\emptyset} = \langle \rangle$
        \item $i=0\ Z^{1} = \langle x_1 \rangle$
    \end{itemize}
}

\subsection{Passo ricorsivo $i > 1$}
{
    \large

    \begin{itemize}
        \item $Z^{(i)} \stackrel{???}{=} Z^{(0)}\ Z^{(1)} \dots Z^{(i-1)}$
        \item $Z^{(i)} = Z^{(?)} + x_i$
    \end{itemize}
}

Non è possibile costruire una LIS in questo modo, essendo che per capire chi è $Z^{(?)}$ \uline{occorre conoscere con cosa termina}, ma ciò non si può sapere essendo la LIS costruita attraverso \textbf{chiamate ricorsive}.

Per risolvere il dilemma, si \uline{introduce un problema ausiliario}/vincolato associato a LIS per ovviare a tale mancanza

\newpage
\section{Formulazione ricorsiva di una $LIS_V$}
\subsection{Istanza}
{
    \large
    $$
        X=X_m=\langle x_1, \dots, x_m \rangle
    $$
}
\subsection{Soluzione}
Una $LIS\ Z$ di $X$ la quale termina con l'ultimo simbolo di $X$ (cioè con $x_m$)

\subsection{Definizione dei coefficienti nel problema $LIS_V$}
$\forall i \in \{1, \dots, m\}\ c_i$ è $|Z^{(i)}|$ cioè la lunghezza di $Z^{(i)}$

\subsection{Caso base}
Con $i=1$ abbiamo che l'istanza è $X_1 =\langle x_1 \rangle$, questo porta ad avere $Z^{(i)}=\langle x_1 \rangle$ e $c_1  = |Z^{(i)}| = 1$

\subsection{Passo ricorsivo}
{
    \large $$ Z^{(i) }= Z^{(k)} + x_i $$ 
}

$Z^(k)$ è uno dei sottoproblemi più piccoli di $i$ che termina con $x_i$.

Si richiede che:
\begin{itemize}
    \item $\mathbf{x_k < x_i}$
    \item $Z^{(k)}$ sia una migliora (più lunga) tra tutte le $Z^{(j)}$ con $j<i$ il cui ultimo simbolo $x_j$ è minore di $x_i$
\end{itemize}

Quindi $Z^{(k)}$ è tale che
{
    \large
    $$ 
        c_k=|Z^{(k)}| = max\{c_j: 1 \leq j \leq \land x_j < x_i\}
    $$
}
siccome $Z^{(i)}=Z^{(j)} + x_i$, $c_i = max\{c_j : 1 \leq i \land x_j < x_i\} + 1$ con $max\ \emptyset=0$

\newpage
\section{Algoritmo Bottom-Up}
\begin{algorithm}
    \caption{LIS-LENGTH($n$)}
    \begin{algorithmic}
        \Procedure{LIS-Length}{$n$}
        
        \State $n=length(X)$

        \State $c[1]=1$ \Comment{Caso base}

        \State $m = c[1]$ \Comment{in m si vuole calcolare il $max\{c_1,\dots,c_n\}$}

        
        \State $ind-m=1$ \Comment{in ind-m si vuole calcolare l'indice $k$ tale che $c_k = max\{c1,\dots,c_n\}$}

        \State $b[1]=0$

        \For{$i=2$ to $n$} \Comment{$\text{occorre calcolare } c_i = max\{c_j:1 \leq j \leq i \land x_j < x_i\} + 1$ con $max\ \emptyset = 0$}

        \State $tmp = 0$
        \State $h = 0$
        
        \For{$j=1$ to $i-1$}
            \If{$x_j < x_i \land c[j]>tmp$}
            \State $tmp=c[j]$
            \State $h=j$
            \EndIf
        \EndFor

        \State $c[i]=tmp+1$
        \State $b[i]=h$

        \If{$c[i]>m$}
        \State $m=c[i]$
        \State $ind-m=i$
        \EndIf

        \EndFor


        \State \Return $c,b,m,ind-m$
        \EndProcedure
    \end{algorithmic}
\end{algorithm}

Complessità in tempo $\Theta(n^2)$ in spazio $\Theta(n)$


\section{Ricostruzione della LIS}
\begin{algorithm}
    \caption{PRINT-LIS}
    \begin{algorithmic}
        \Procedure{Print-LIS}{$ind-m,b$}
        \If{$b[ind-m]=0$}
            \State print($x_i$)

            \Else
                \State Print-LIS($b[i]$)
                \State print($x_i$)
        \EndIf
        \EndProcedure
    \end{algorithmic}
\end{algorithm}


\newpage
\section{Dimostrazione: Proprietà Sottostruttura Ottima di una $LIS_V$}
\subsection{Teorema}
\begin{itemize}
    \item Data $X= \langle x_1, \dots, x_n \rangle$
    \item Sia $i\in \{2,\dots,n\}$
    \item Sia $Z^{(i)}$ una $LIS_V$ di $X_i$
    \item Sia $Z^{(i)}$ con $j \in \{1,\dots,i-1\}$ una qualuque $LIS_V$ di $X_j$
    \begin{itemize}
        \item Quindi $Z^{(i)}$ termina con $x_i$ e $Z^{(j)}$ termina con $x_j$
    \end{itemize}
\end{itemize}

Vale che:
{
    \large
    $$Z^{(i)} = Z^* | x_i$$
}

dove $Z^*$ è una più lunga tra le sequenze $Z^{j}$ il cui ultimo simbolo $x_j$ è $<x_i$ 

\subsection{Dimostrazione}
Trattiamo solo il caso in cui $Z^* \not = \langle \rangle$.

\textbf{Se per assurdo} $Z^*|x_i$ non è soluzione del sottoproblema i-esimo per nessuna tra le possibili sequenze $Z^*$ allora 
$\exists Z' : Z^{(i)} = Z'|x_i$ vale quindi che:
\begin{itemize}
    \item $Z'$ è crescente
    \item $Z'$ termina con un simbolo $<x_i$
    \item $|Z'|>|Z^*|$ termina con $x_j$, che è uno tra $x_{j_1},\dots,x_{j_e}$
\end{itemize}

Quindi considerando $Z^* = Z^{(j)}$ si ha che $Z^*$ non è una soluzione del sottoproblema $j'$, il che contraddice l'ipotesi